\chapter{Semiconductor Physics \& Solid State Devices}

\section{Classical Drude-Lorentz Free Electron Theory}
In 1900, Paul Drude treated valence electrons in metals as a classical gas of non-interacting free particles moving among fixed positive ionic cores. Under an applied electric field $\vec{E}$, electrons accelerate according to $m\frac{d\vec{v}}{dt} = -e\vec{E} - \frac{m\vec{v}}{\tau}$.
In steady state, electrons achieve a constant \textbf{Drift Velocity} $v_d$:
\begin{equation}
v_d = \mu E = \frac{e\tau}{m} E
\end{equation}
where $\tau$ is the mean relaxation time between collisions and $\mu = \frac{e\tau}{m}$ is electron mobility.
The current density is $J = n e v_d = \sigma E$, giving \textbf{Electrical Conductivity}:
\begin{equation}
\sigma = \frac{n e^2 \tau}{m}
\end{equation}
Drude theory successfully derived the \textbf{Wiedemann-Franz Law}, stating that the ratio of thermal conductivity $K$ to electrical conductivity $\sigma$ is directly proportional to absolute temperature $T$:
\begin{equation}
\frac{K}{\sigma T} = L = \frac{\pi^2 k_B^2}{3 e^2} \approx 2.44 \times 10^{-8}\text{ W}\cdot\Omega/\text{K}^2
\end{equation}
where $L$ is the universal \textbf{Lorentz Number}.

\subsection{Limitations of Classical Free Electron Theory}
Despite these successes, classical Drude theory suffered severe failures:
\begin{enumerate}[itemsep=2pt]
    \item It predicted an electronic heat capacity $C_{v,\text{el}} = \frac{3}{2} R$, whereas experimental values at room temperature are $100\times$ smaller ($C_{v,\text{el}} \propto T$).
    \item It could not explain why materials with similar electron densities act as conductors, semiconductors, or insulators.
    \item It could not explain positive Hall coefficients observed in metals such as Zinc, Cadmium, and Beryllium.
\end{enumerate}

\section{Quantum Band Theory and the Kronig-Penney Model}
Felix Bloch (1928) demonstrated that an electron moves inside a strictly periodic crystal lattice potential $V(x+a) = V(x)$. According to \textbf{Bloch's Theorem}, the stationary wavefunction is a modulated plane wave:
\begin{equation}
\psi_k(x) = e^{ikx} u_k(x), \quad \text{where } u_k(x+a) = u_k(x)
\end{equation}

\subsection{The Kronig-Penney Potential Model}
Kronig and Penney (1931) modeled the 1D crystal as a periodic array of rectangular potential wells of width $a$, barrier width $b$, and barrier height $V_0$.
Solving the Schrödinger equation in the well ($0 < x < a$) and barrier ($-b < x < 0$) and matching boundary conditions at the interfaces yields the transcendental dispersion equation:
\begin{equation}
P\frac{\sin(\alpha a)}{\alpha a} + \cos(\alpha a) = \cos(ka)
\end{equation}
where $P = \lim_{b \to 0, V_0 \to \infty} \left( \frac{m V_0 b a}{\hbar^2} \right)$ is the \textbf{barrier strength / potential binding parameter}, and $\alpha = \frac{\sqrt{2mE}}{\hbar}$.

\begin{conceptbox}{Key Physical Insights from the Kronig-Penney Model}
Because $|\cos(ka)| \le 1$, real wave vector solutions $k$ exist only when the function $f(\alpha a) = P\frac{\sin(\alpha a)}{\alpha a} + \cos(\alpha a)$ falls in the range $[-1, +1]$.
\begin{enumerate}[itemsep=2pt]
    \item The electronic energy spectrum splits into alternating \textbf{Allowed Energy Bands} (where $|f(\alpha a)| \le 1$) and \textbf{Forbidden Energy Gaps ($E_g$)} (where $|f(\alpha a)| > 1$).
    \item As electron energy increases ($\alpha a$ increases), the allowed bands widen and the forbidden energy bandgaps shrink.
    \item If $P \to 0$ (free electron limit), $\cos(\alpha a) = \cos(ka) \implies \alpha = k \implies E = \frac{\hbar^2 k^2}{2m}$ (continuous parabolic energy spectrum without bandgaps).
    \item If $P \to \infty$ (isolated atom limit), $\sin(\alpha a) = 0 \implies \alpha a = n\pi \implies E_n = \frac{n^2 \pi^2 \hbar^2}{2ma^2}$ (discrete quantized atomic levels).
\end{enumerate}
\end{conceptbox}

\section{Effective Mass and the Concept of Holes}
Inside a crystal lattice, an electron experiences both external electric fields and intense internal periodic crystal forces. By applying wave packet dynamics ($v_g = \frac{1}{\hbar}\frac{dE}{dk}$), acceleration is $a = \frac{dv_g}{dt} = \frac{1}{\hbar}\frac{d^2E}{dk^2}\frac{dk}{dt}$.
Using $\hbar \frac{dk}{dt} = F_{\text{ext}} = -eE$, we obtain $a = \frac{1}{\hbar^2}\frac{d^2E}{dk^2} F_{\text{ext}}$.
Comparing this with Newton's second law ($F = m^* a$) defines the \textbf{Effective Mass} $m^*$:
\begin{equation}
m^* \equiv \frac{\hbar^2}{\frac{d^2 E}{dk^2}}
\end{equation}
\begin{itemize}[itemsep=2pt]
    \item Near the bottom of a conduction band ($\frac{d^2E}{dk^2} > 0$): $m^* > 0$. The electron behaves as a positively massed particle accelerated in the field direction.
    \item Near the top of a valence band ($\frac{d^2E}{dk^2} < 0$): $m^* < 0$. An electron with negative effective mass accelerates opposite to the normal electronic trajectory, behaving mathematically as a positively charged particle: a \textbf{Hole} ($q = +e, m_h^* > 0$).
\end{itemize}

\section{Fermi-Dirac Carrier Statistics}
Electrons are indistinguishable fermions obeying the Pauli exclusion principle. The probability that an available quantum state at energy $E$ is occupied at temperature $T$ is given by the \textbf{Fermi-Dirac Distribution Function}:
\begin{equation}
f(E) = \frac{1}{1 + \exp\left(\frac{E - E_F}{k_B T}\right)}
\end{equation}
where $E_F$ is the \textbf{Fermi Energy Level}.
\begin{itemize}[itemsep=2pt]
    \item \textbf{At Absolute Zero ($T = 0\text{ K}$):}
    \[
    f(E) = \begin{cases} 1 & \text{for } E < E_F \quad \text{(all states fully occupied)} \\ 0 & \text{for } E > E_F \quad \text{(all states completely empty)} \end{cases}
    \]
    \item \textbf{At Finite Temperature ($T > 0\text{ K}$):} At $E = E_F$, $f(E_F) = \frac{1}{1 + e^0} = 0.5$ ($50\%$ occupation probability).
\end{itemize}

\subsection{Intrinsic and Extrinsic Carrier Densities}
The density of electrons in the conduction band and holes in the valence band are:
\begin{equation}
n = N_c \exp\left(-\frac{E_c - E_F}{k_B T}\right), \quad p = N_v \exp\left(-\frac{E_F - E_v}{k_B T}\right)
\end{equation}
where $N_c = 2\left(\frac{2\pi m_e^* k_B T}{h^2}\right)^{3/2}$ and $N_v = 2\left(\frac{2\pi m_h^* k_B T}{h^2}\right)^{3/2}$ are the effective density of states.

For an \textbf{intrinsic semiconductor} ($n = p = n_i$):
\begin{equation}
n_i = \sqrt{N_c N_v} e^{-E_g / (2k_B T)}
\end{equation}
\begin{equation}
E_F^{\text{intrinsic}} = \frac{E_c + E_v}{2} + \frac{3}{4}k_B T \ln\left(\frac{m_h^*}{m_e^*}\right) \approx \frac{E_c + E_v}{2} \quad (\text{Exact middle of bandgap})
\end{equation}
\textbf{Mass Action Law:} For any semiconductor under thermal equilibrium:
\begin{equation}
n \cdot p = n_i^2 = N_c N_v e^{-E_g / k_B T}
\end{equation}


\section{Density of States in 3D Semiconductor Bands}
To calculate carrier densities, we determine the number of quantum electron states per unit energy range per unit volume: the \textbf{Density of States} $g(E)$.
Applying periodic boundary conditions in 3D $k$-space ($V = L^3$):
\begin{equation}
g(E)\,dE = 2 \times \frac{4\pi k^2 \, dk}{(2\pi/L)^3 \cdot L^3} = \frac{k^2}{\pi^2}\frac{dk}{dE} \, dE
\end{equation}
Using the parabolic conduction band dispersion $E - E_c = \frac{\hbar^2 k^2}{2m_e^*}$:
\begin{equation}
g_c(E) = \frac{1}{2\pi^2}\left(\frac{2m_e^*}{\hbar^2}\right)^{3/2}\sqrt{E - E_c} \quad (E \ge E_c)
\end{equation}
For the valence band:
\begin{equation}
g_v(E) = \frac{1}{2\pi^2}\left(\frac{2m_h^*}{\hbar^2}\right)^{3/2}\sqrt{E_v - E} \quad (E \le E_v)
\end{equation}

\section{Direct vs Indirect Bandgap Semiconductors}
\begin{itemize}[leftmargin=*,itemsep=3pt]
    \item \textbf{Direct Bandgap Semiconductors (e.g., GaAs, InP, GaN):} The conduction band minimum and valence band maximum occur at the identical crystal wave vector ($k = 0$). An electron can recombine with a hole via direct emission of a photon without lattice phonon assistance:
    \begin{equation}
    E_{\text{photon}} = h\nu \approx E_g, \quad \Delta k \approx 0
    \end{equation}
    Radiative recombination rate is extremely high ($\tau \sim 1\text{--}10\text{ ns}$), making them ideal for LEDs, laser diodes, and high-efficiency solar cells.
    \item \textbf{Indirect Bandgap Semiconductors (e.g., Si, Ge, AlAs):} The conduction band minimum is offset in $k$-space from the valence band maximum ($k_c \ne 0$). Recombination requires both a photon (for energy conservation) and a lattice vibration / phonon (for momentum conservation):
    \begin{equation}
    E_e - E_h = h\nu \pm \hbar\Omega_{\text{phonon}}, \quad k_e - k_h = \vec{q}_{\text{phonon}}
    \end{equation}
    This three-particle interaction is vastly slower ($\tau \sim 1\text{--}100\text{ }\mu\text{s}$), and energy is lost as thermal heat. Silicon is therefore unsuitable for light-emitting lasers without nanostructuring.
\end{itemize}

\section{The Hall Effect and Semiconductor Characterization}
Discovered by Edwin Hall in 1879. When a current $I$ flows along the $x$-axis of a semiconductor slab of width $w$ and thickness $t$ placed in a transverse magnetic field $\vec{B} = B\hat{k}$, the Lorentz force $\vec{F}_B = q(\vec{v}_d \times \vec{B})$ deflects charge carriers along the $y$-axis.
Accumulated charges create a transverse \textbf{Hall Electric Field} $\vec{E}_H = E_H \hat{j}$ that opposes further deflection:
\begin{equation}
q E_H = q v_d B \implies E_H = v_d B
\end{equation}
Using current density $J = \frac{I}{w t} = n q v_d \implies v_d = \frac{I}{n q w t}$, the measured \textbf{Hall Voltage} $V_H$ is:
\begin{equation}
V_H = E_H w = v_d B w = \left(\frac{I}{n q w t}\right) B w = \frac{B I}{n q t} = \frac{R_H B I}{t}
\end{equation}
where the \textbf{Hall Coefficient} is defined as:
\begin{equation}
R_H \equiv \frac{1}{n q} = \begin{cases} -\frac{1}{n e} & \text{(N-type semiconductor, electron majority)} \\ +\frac{1}{p e} & \text{(P-type semiconductor, hole majority)} \end{cases}
\end{equation}

\begin{takeawaybox}
\textbf{Crucial Engineering Applications of the Hall Effect:}
\begin{enumerate}[itemsep=2pt]
    \item \textbf{Carrier Type Determination:} The polarity of $V_H$ immediately identifies N-type vs P-type material.
    \item \textbf{Carrier Concentration Measurement:} Carrier density is directly computed as $n = \frac{1}{e |R_H|}$.
    \item \textbf{Carrier Mobility Calculation:} Hall mobility is calculated via $\mu_H = |R_H| \sigma$.
    \item \textbf{Magnetic Field Sensing:} Used in contactless automotive wheel-speed sensors, brushless DC motor commutators, and digital compasses.
\end{enumerate}
\end{takeawaybox}

\section{Photovoltaic Solar Cells}
A solar cell is a large-area semiconductor p-n junction that converts solar radiation directly into electrical power through the photovoltaic effect:
\begin{enumerate}[leftmargin=*,itemsep=2pt]
    \item \textbf{Electron-Hole Generation:} Incident photons with energy $h\nu \ge E_g$ excite electrons across the bandgap, creating electron-hole pairs.
    \item \textbf{Carrier Collection:} The built-in electric field in the depletion zone sweeps electrons into the N-region and holes into the P-region before they can recombine.
    \item \textbf{Electrical Characteristics:}
    \begin{equation}
    I = I_{\text{sc}} - I_0 \left( e^{qV / \eta k_B T} - 1 \right)
    \end{equation}
    \item \textbf{Fill Factor (FF) and Conversion Efficiency ($\eta$):}
    \begin{equation}
    \text{FF} \equiv \frac{V_{\text{mp}} I_{\text{mp}}}{V_{\text{oc}} I_{\text{sc}}}, \quad \eta = \frac{P_{\text{max}}}{P_{\text{in}}} = \frac{\text{FF} \cdot V_{\text{oc}} \cdot I_{\text{sc}}}{P_{\text{in}} \cdot A_{\text{cell}}}
    \end{equation}
\end{enumerate}

\section{Worked Numerical Examples}

\begin{examplebox}{Worked Example 5.1: Hall Effect Carrier Concentration and Mobility}
\textbf{Problem:} A rectangular Silicon strip of thickness $t = 0.5\text{ mm}$ and width $w = 2\text{ mm}$ carries a current $I = 10\text{ mA}$ in a magnetic field $B = 0.5\text{ Tesla}$. The measured Hall voltage is $V_H = -2.5\text{ mV}$. The electrical conductivity of the sample is $\sigma = 100\text{ S/m}$. Calculate: (a) The Hall coefficient $R_H$, (b) The carrier type and density $n$, and (c) The carrier mobility $\mu_H$.\\
\textbf{Solution:}\\
(a) Hall coefficient:
\[
R_H = \frac{V_H t}{B I} = \frac{(-2.5 \times 10^{-3}\text{ V}) \times (0.5 \times 10^{-3}\text{ m})}{(0.5\text{ T}) \times (10 \times 10^{-3}\text{ A})} = \frac{-1.25 \times 10^{-6}}{2.5 \times 10^{-3}} = -5.0 \times 10^{-4}\text{ m}^3/\text{C}
\]
(b) Since $R_H < 0$, the sample is an \textbf{N-type semiconductor} (electron conduction).\\
Electron concentration:
\[
n = \frac{1}{e |R_H|} = \frac{1}{(1.602 \times 10^{-19}\text{ C}) \times (5.0 \times 10^{-4}\text{ m}^3/\text{C})} = \frac{1}{8.01 \times 10^{-23}} = 1.25 \times 10^{22}\text{ m}^{-3}
\]
(c) Hall mobility:
\[
\mu_H = |R_H| \sigma = (5.0 \times 10^{-4}\text{ m}^3/\text{C}) \times (100\text{ S/m}) = 0.05\text{ m}^2/(\text{V}\cdot\text{s}) = 500\text{ cm}^2/(\text{V}\cdot\text{s})
\]
\end{examplebox}

\section{Review Questions}
\begin{enumerate}[leftmargin=*,itemsep=3pt]
    \item Explain the Kronig-Penney model and derive how periodic lattice potentials lead to allowed bands and forbidden bandgaps.
    \item Define effective mass ($m^*$). Why is the effective mass negative near the top of a valence band?
    \item State Fermi-Dirac distribution. Plot $f(E)$ vs $E$ at $T = 0\text{ K}$ and $T > 0\text{ K}$.
    \item Derive the expression for Hall voltage $V_H$ and list four major engineering applications of the Hall effect.
\end{enumerate}
